\section{Interpretation and extensions}

The rate characterization in \cref{obj:thm:uniform-frontier} describes the information available when assignment records the latent cutoff side, the latent score has known bounded support, and the Gaussian measurement-error distribution is calibrated. Recorded assignment separates the observations into the two latent-score arms. Within each arm, the cutoff is an endpoint, and the contaminated score supplies information about distance from that endpoint. These features give the endpoint weights in \cref{obj:def:endpoint-kernel} their statistical interpretation.

Known Gaussian calibration supplies the inverse-heat transform in \cref{obj:def:inverse-heat}. The finite-sample guarantees in \cref{obj:thm:finite-certificate} combine this correction with the global density and conditional-mean restrictions defining \cref{obj:def:model}. The density bounds control the mass entering the arm ratios, while the full-interval Hölder restrictions control their approximation error. Together, these conditions support observable estimation and honest inference uniformly over the benchmark, with matching minimax orders across the stated noise scales.

\subsection{Converse transfer to a larger Gaussian class}

The benchmark also supplies lower bounds for a larger class specified through local density and conditional-mean conditions. To express this implication precisely, the following definition fixes the larger class, its target, and its decision criteria immediately before the embedding result.



\begin{propositionv}{prop:published-class-converse-transfer}[Gaussian class converse transfer]
For \(\sigma\in[0,1]\), let \(\mathcal G(\sigma)\) consist of probability laws \(Q\) on \(\mathbb R\times\{0,1\}^2\times\mathbb R\), with coordinates \((Z,V^0,V^1,E)\), a score density \(a\), and conditional-mean versions
\[
v_d(z)=\mathbb E_Q[V^d\mid Z=z],\qquad d\in\{0,1\}.
\]
Define the assignment, observed outcome, and observed score by
\[
A=\mathbf 1\{Z\ge0\},\qquad
B=AV^1+(1-A)V^0,\qquad
R=Z+\sigma E.
\]
Membership requires the following conditions:
\begin{itemize}
\item \textbf{(Local density.)} There exists \(r>0\) such that \(a\) is continuous on \((-r,r)\) and \(a(z)>0\) for every \(z\in(-r,r)\).
\item \textbf{(Conditional means.)} Each \(v_d\) is a conditional-mean version continuous on an open neighborhood of zero, with \(0\le v_d(0)\le1\).
\item \textbf{(Gaussian error.)} \(E\sim N(0,1)\).
\item \textbf{(Joint independence.)} \(E\perp(Z,B)\).
\end{itemize}
The target and randomized experiment are
\[
\vartheta(Q)=v_1(0)-v_0(0),\qquad
\bigl(Q\circ(R,A,B)^{-1}\bigr)^{\otimes n}
 \otimes\operatorname{Unif}[0,1].
\]
Let \(\mathcal R_{\mathcal G}(n,\sigma)\) denote the infimum, over measurable estimates based on this experiment, of
\[
\sup_{Q\in\mathcal G(\sigma)}
\mathbb E_Q\!\left[\,|T-\vartheta(Q)|\,\right].
\]
Let \(\mathcal L_{\mathcal G}(n,\sigma)\) denote the infimum of
\[
\sup_{Q\in\mathcal G(\sigma)}
\mathbb E_Q[\upsilon-\ell]
\]
over measurable connected closed intervals \([\ell,\upsilon]\subseteq[-1,1]\) based on the same experiment and satisfying
\[
\inf_{Q\in\mathcal G(\sigma)}
Q\{\vartheta(Q)\in[\ell,\upsilon]\}\ge0.9.
\]
These expectations and coverage probabilities include the independent uniform procedure seed.

Fix \(\beta\in(0,1]\). For every \(\sigma\in[0,1]\) and every \(P\in\mathcal M_\beta(\sigma)\), the law \(P\), equipped with the density \(a=f(P,\cdot)\) and versions \(v_d=\mu_d(P,\cdot)\), belongs to \(\mathcal G(\sigma)\). This embedding preserves the target,
\[
\vartheta(Q)=\theta(P)
=\mu_1(P,0)-\mu_0(P,0),
\]
the observed law,
\[
Q\circ(R,A,B)^{-1}=P_{\mathrm{obs}},
\]
and, for every integer \(n\ge0\), the entire randomized experiment,
\[
\bigl(Q\circ(R,A,B)^{-1}\bigr)^{\otimes n}
 \otimes\operatorname{Unif}[0,1]
=\mathbb P_{P,n}.
\]

There exists \(c>0\), depending on \(\beta\), such that for every integer \(n\ge2\) and every \(\sigma\in[0,1]\),
\[
c\,r_\beta(n,\sigma)
\le R_\beta(n,\sigma)
\le \mathcal R_{\mathcal G}(n,\sigma),
\qquad
c\,r_\beta(n,\sigma)
\le \mathcal L_\beta(n,\sigma)
\le \mathcal L_{\mathcal G}(n,\sigma),
\]
where \(r_\beta(n,\sigma)\) is the benchmark frontier rate, and \(R_\beta(n,\sigma)\) and \(\mathcal L_\beta(n,\sigma)\) are the corresponding minimax absolute risk and minimax expected length under uniform \(0.9\) coverage on \(\mathcal M_\beta(\sigma)\).

For every \(\sigma\in(0,1]\) and every \(Q\in\mathcal G(\sigma)\), the locally defined binary-response Pei--Shen conditions hold: the standardized error is independent jointly of the latent score and observed outcome (Assumption 1Y), the score density is strictly positive on an open neighborhood of zero (Assumption 2C), and the scaled error satisfies
\[
\sigma E\sim N(0,\sigma^2)
\]
(Assumption 7).
\end{propositionv}

The intuition behind \cref{obj:prop:published-class-converse-transfer} is inclusion of statistical decision problems with exactly the same target and data distribution on the embedded benchmark. Any estimator evaluated over the larger class faces every benchmark law. Likewise, an interval honest over the larger class satisfies the benchmark coverage requirement. Preserving the independent procedure seed extends both comparisons to randomized procedures. The benchmark lower bounds therefore remain necessary precision bounds for the larger class.

The local positivity and continuity requirements in the larger Gaussian class of \cref{obj:prop:published-class-converse-transfer} give the cutoff contrast a well-defined meaning, while joint independence specifies a common Gaussian contamination mechanism for the latent score and observed outcome. For positive noise scales, that proposition connects these restrictions to Assumptions~1Y, 2C, and 7 of \citet{PeiShen2017DevilTails}. The resulting interpretation is an embedding-based converse for the larger Gaussian class. Matching attainment and finite-sample honesty have the exact benchmark in \cref{obj:def:model} as their domain, as established by \cref{obj:thm:finite-certificate,obj:thm:uniform-frontier}.

\subsection{A prospective eligibility-effect benchmark}

Following the Medicaid/CHIP setting in \citet[Chapter~2, Section~2.6]{Dong2018Thesis}, the proposed eligibility application would be organized around an eligibility cutoff.
% [Dong dissertation, Chapter 2, Section 2.6](https://etheses.lse.ac.uk/3756/1/Dong__essays-in-microeconometrics.pdf)
The proposed design would link administrative eligibility assignment to a contaminated score and a prespecified binary outcome. The score would be oriented so that eligibility corresponds to its nonnegative latent side. Potential outcomes would be indexed by eligibility assignment, giving the cutoff contrast the interpretation of an eligibility effect on the chosen outcome.

Applying the benchmark would require an administrative record that identifies the latent cutoff side, together with externally justified Gaussian calibration of the score error and its known scale. The proposed error mechanism would satisfy joint independence from the latent score and both potential outcomes. A prespecified population would have a known bounded latent range, normalized to \([-1,1]\) with the cutoff at zero, and observations would follow the independent sampling experiment in \cref{obj:def:model}.

The remaining proposed design conditions are substantive global restrictions: a continuous latent density between \(1/4\) and \(3/4\) throughout the normalized interval, continuous binary potential-outcome means between the same bounds, and full-interval Hölder seminorms at most two for a specified exponent in \((0,1]\). These are the restrictions in \cref{obj:ass:density-bounds,obj:ass:mean-bounds,obj:ass:mean-holder}. Under these maintained conditions, \cref{obj:thm:finite-certificate} supplies an observable interval with uniform finite-sample coverage, and \cref{obj:thm:uniform-frontier} characterizes its attainable precision. The application thus provides a prospective design benchmark conditional on administrative linkage, calibration, support, and global regularity.

\subsection{Limitations and future work}

The analysis treats the Gaussian error distribution and its scale as known. Estimated calibration would introduce uncertainty into both the inverse-heat weights and the public interval radius. Extending the coverage guarantee to that setting requires an experiment that includes the calibration data and conditions controlling its estimation error.

Smoothness is also a public input. Adaptation to an unknown exponent would require procedures and coverage arguments that account for selection across smoothness classes. Empirical assessment of benchmark membership presents a separate challenge: the support, density bounds, and global potential-outcome mean restrictions concern latent quantities. Administrative and validation data could inform their plausibility, but the prospective eligibility design remains conditional on these maintained restrictions.

Additional outcome models would broaden the application beyond binary responses. Such extensions require suitable outcome-tail or moment conditions and corresponding changes to the finite-sample second-moment bounds. The present results establish estimation and honest inference for the binary-outcome benchmark.

Finally, the focused comparison in \cref{sec:related-work} separates three antecedents: boundary kernels estimate a truncated latent density, deconvolution regression uses a numerator--denominator ratio at fixed error law, and convolution linear-functional theory gives quadratic-risk bounds for individual functionals. None of these cited results supplies the joint criterion used here: a recorded-arm endpoint contrast with absolute-risk and honest-length guarantees uniform over the benchmark noise scales. The comparison with \citet{ZhangKarunamuni2009Boundary} is qualitative, covering the boundary-density target, error classes, and pointwise mean squared error criterion described in the available primary-source material. The operational size of the explicit finite-sample certificate at applied sample sizes is likewise left for a separate calibrated numerical study; the present paper establishes its validity and optimal order without reporting a design-specific grid.
